%% toc: P15 166 166.6
\section{Maximum-entropy inference under declared constraints}
\label{sec:annex-p15-w02b-p15-043-maxent}

Throughout, the notation is that of
\cref{def:p15-w02b-p15-043-problem,def:p15-w02b-p15-043-partition}.
The function \(\varphi(t)\defeq t\ln t-t+1\) for \(t>0\), with
\(\varphi(0)\defeq1\), is continuous on \([0,\infty)\). Since
\(\varphi'(t)=\ln t\) is negative on \((0,1)\) and positive on
\((1,\infty)\), and \(\varphi(1)=0\), one has \(\varphi\geq0\) with
equality only at \(t=1\). For \(x\geq0\) and \(y>0\) this gives
\begin{equation*}
 x\ln\frac xy-(x-y)=y\,\varphi(x/y)\geq0,
 \qquad\text{with equality only if }x=y. \tag{\(\ast\)}
\end{equation*}

\begin{annexitem}{annex:p15-w02b-p15-043-gibbs}{thm:p15-w02b-p15-043-gibbs}{The Gibbs identity and the dual bound}
\begin{proof}
(1) Let \(r=p/q\), defined \(\nu\)-almost everywhere. By \((\ast)\), almost
everywhere
\[
 p\ln\frac pq=q\,\varphi(r)+(p-q),
\]
both sides being \(0\) where \(p=0\). The first term is nonnegative and the
second is \(\nu\)-integrable with integral \(1-1=0\). Hence the negative part
of \(p\ln(p/q)\) is at most \(|p-q|\), which is integrable, so the integral is
defined, and
\(D(P\Vert Q)=\int q\,\varphi(r)\,d\nu\in[0,\infty]\). It vanishes exactly
when \(q\varphi(r)=0\) almost everywhere; since \(q>0\) almost everywhere this
means \(\varphi(r)=0\), that is \(r=1\), almost everywhere. Thus
\(D(P\Vert Q)=0\) if and only if \(p=q\) almost everywhere, that is \(P=Q\).

(2) The density \(q_\lambda=\exp(-\lambda_0-\lambda\cdot f)\) is positive
everywhere. Since \(|\ln q_\lambda|\leq|\lambda_0|+\sum_k|\lambda_k||f_k|\)
and each \(f_k\) is \(P\)-integrable, \(\ln q_\lambda\in L^1(P)\) and
\[
 \int p\ln q_\lambda\,d\nu=E_P[\ln q_\lambda]=-\lambda_0-\sum_k\lambda_kc_k .
\]
Where \(p>0\), \(p\ln p=p\ln(p/q_\lambda)+p\ln q_\lambda\); where \(p=0\)
both sides are \(0\). The first summand has integrable negative part by (1)
and the second is integrable. Hence \((p\ln p)^-\) is integrable, the integral
of \(p\ln p\) is defined in \((-\infty,\infty]\), and
\[
 \int p\ln p\,d\nu=D(P\Vert Q_\lambda)-\lambda_0-\sum_k\lambda_kc_k .
\]
Changing sign gives the stated identity, and \(H_\nu(P)\leq\Psi_c(\lambda)\)
because \(D\geq0\).

(3) If \(\mathcal F_\nu(c)\) is empty, \(v_\nu(c)=-\infty\). If
\(\Lambda_\nu\) is empty, the infimum is \(+\infty\). Otherwise every feasible
\(P\) has defined entropy by (2), and \(H_\nu(P)\leq\Psi_c(\lambda)\) for
every \(\lambda\in\Lambda_\nu\). Taking the supremum over \(P\) and then the
infimum over \(\lambda\) gives the bound. No maximizer was used.
\end{proof}
\end{annexitem}

\begin{annexitem}{annex:p15-w02b-p15-043-optimality}{thm:p15-w02b-p15-043-optimality}{A feasible exponential-family member is the unique maximizer}
\begin{proof}
By hypothesis \(Q_\lambda\ll\nu\), each \(f_k\) is \(Q_\lambda\)-integrable and
\(E_{Q_\lambda}[f_k]=c_k\); so \(Q_\lambda\in\mathcal F_\nu(c)\). Part (2) of
\cref{thm:p15-w02b-p15-043-gibbs} with \(P=Q_\lambda\) and
\(D(Q_\lambda\Vert Q_\lambda)=0\) gives
\(H_\nu(Q_\lambda)=\Psi_c(\lambda)=\lambda_0+\sum_k\lambda_kc_k\).
For any feasible \(P\), the same part gives
\(H_\nu(P)=\Psi_c(\lambda)-D(P\Vert Q_\lambda)\leq H_\nu(Q_\lambda)\), and
part (1) shows that equality forces \(P=Q_\lambda\). Hence
\(v_\nu(c)=H_\nu(Q_\lambda)\) and \(Q_\lambda\) is the only maximizer.
\end{proof}
\end{annexitem}

\begin{annexitem}{annex:p15-w02b-p15-043-no-constraint}{thm:p15-w02b-p15-043-no-constraint}{No constraints: normalized reference measure or unbounded entropy}
\begin{proof}
With \(m=0\) the constraints are vacuous and \(\mathbb R^0\) has one point,
at which \(Z_\nu=\nu(\mathcal X)\). If \(\nu(\mathcal X)<\infty\), that point
lies in \(\Lambda_\nu\), \(\lambda_0=\ln\nu(\mathcal X)\), and \(Q\) has
constant density \(1/\nu(\mathcal X)\), so \(Q=\nu/\nu(\mathcal X)\).
\Cref{thm:p15-w02b-p15-043-optimality} gives uniqueness and the value
\(\lambda_0=\ln\nu(\mathcal X)\). In the finite case
\(Q\{x\}=w(x)/\sum_yw(y)\); with \(w\equiv1\) on \(n\) points this is
\(1/n\), and the value is \(\ln n\).

If \(\nu(\mathcal X)=\infty\), \(\sigma\)-finiteness gives sets
\(\mathcal X_1,\mathcal X_2,\ldots\) of finite measure with union
\(\mathcal X\). Put \(A_j=\mathcal X_1\cup\cdots\cup\mathcal X_j\). Then
\(\nu(A_j)<\infty\) and, by continuity from below,
\(\nu(A_j)\uparrow\infty\). For \(j\) with \(\nu(A_j)>0\), the probability
\(U_j\) with density \(\mathbf 1_{A_j}/\nu(A_j)\) is feasible and
\[
 H_\nu(U_j)=-\int_{A_j}\frac1{\nu(A_j)}\ln\frac1{\nu(A_j)}\,d\nu=\ln\nu(A_j)\to\infty .
\]
So \(v_\nu=+\infty\), and no feasible distribution of finite entropy attains it.
\end{proof}
\end{annexitem}

\begin{annexitem}{annex:p15-w02b-p15-043-uniqueness}{thm:p15-w02b-p15-043-uniqueness}{Strict concavity and uniqueness}
\begin{proof}
\emph{Convexity.} For feasible \(P,P'\) and \(t\in[0,1]\), the measure
\(tP+(1-t)P'\) is a probability measure with density \(tp+(1-t)p'\) relative
to \(\nu\); each \(|f_k|(tp+(1-t)p')\) is integrable, and by linearity its
\(f_k\)-expectation is \(tc_k+(1-t)c_k=c_k\).

\emph{Strict concavity.} Let \(m=\tfrac12(p_1+p_2)\), the density of \(M\).
Define \(g=\tfrac12m\bigl(\varphi(p_1/m)+\varphi(p_2/m)\bigr)\) where
\(m>0\) and \(g=0\) where \(m=0\). Then \(g\geq0\). Where \(m>0\), the
definition of \(\varphi\) and \(p_1+p_2=2m\) give
\[
 g=\tfrac12\Bigl(p_1\ln\frac{p_1}m+p_2\ln\frac{p_2}m\Bigr)
  =\tfrac12(p_1\ln p_1+p_2\ln p_2)-m\ln m ,
\]
and this also holds where \(m=0\), since then \(p_1=p_2=0\). Because
\(p_i/m\leq2\), \(g\leq m\ln2\), so \(g\) is integrable. The entropies of
\(P_1,P_2\) are finite, so \(p_i\ln p_i\in L^1(\nu)\). Therefore
\(-m\ln m=g-\tfrac12(p_1\ln p_1+p_2\ln p_2)\) is integrable and
\[
 H_\nu(M)=\tfrac12\bigl(H_\nu(P_1)+H_\nu(P_2)\bigr)+\int g\,d\nu
        \geq\tfrac12\bigl(H_\nu(P_1)+H_\nu(P_2)\bigr).
\]
Equality forces \(g=0\) almost everywhere. Where \(m>0\), this means
\(\varphi(p_1/m)=\varphi(p_2/m)=0\), so \(p_1=p_2=m\); where \(m=0\),
\(p_1=p_2=0\). Thus \(P_1=P_2\).

\emph{Uniqueness.} If \(P_1,P_2\) are maximizers with finite value \(v\), then
\(M\) is feasible with defined entropy, so \(H_\nu(M)\leq v\). The display gives
\(H_\nu(M)=v+\int g\,d\nu\), so \(\int g\,d\nu=0\) and \(P_1=P_2\).
\end{proof}
\end{annexitem}

\begin{annexitem}{annex:p15-w02b-p15-043-partition}{thm:p15-w02b-p15-043-partition}{Properties of the partition function}
\begin{proof}
(1) Let \(\lambda,\lambda'\in\Lambda_\nu\) and \(\theta\in(0,1)\). Since
\[
 e^{-(\theta\lambda+(1-\theta)\lambda')\cdot f}
 =(e^{-\lambda\cdot f})^\theta(e^{-\lambda'\cdot f})^{1-\theta},
\]
H\"older's inequality with exponents \(1/\theta\) and \(1/(1-\theta)\) gives
\[
 Z_\nu(\theta\lambda+(1-\theta)\lambda')\leq Z_\nu(\lambda)^\theta Z_\nu(\lambda')^{1-\theta}<\infty .
\]
So \(\Lambda_\nu\) is convex, and taking logarithms shows that \(\ln Z_\nu\) is
convex on it. If \(\lambda_j\to\lambda\), then
\(e^{-\lambda_j\cdot f(x)}\to e^{-\lambda\cdot f(x)}\) for every \(x\), and
Fatou's lemma gives \(Z_\nu(\lambda)\leq\liminf_jZ_\nu(\lambda_j)\). Also
\(Z_\nu>0\) everywhere, since the integrand is positive and \(\nu\neq0\).

(2) Fix an interior point \(\lambda\) and \(u\neq0\); put \(s=u\cdot f\).
Choose \(\varepsilon>0\) with \(\lambda\pm2\varepsilon u\in\Lambda_\nu\).
For real \(y\) and \(j\in\{0,1,2,3\}\), the series of \(e^{\varepsilon|y|}\)
gives \(|y|^j\leq j!\,\varepsilon^{-j}e^{\varepsilon|y|}\). Hence, for
\(|t|\leq\varepsilon\) and \(j\leq2\),
\[
 |s|^je^{-(\lambda+tu)\cdot f}\leq|s|^je^{\varepsilon|s|}e^{-\lambda\cdot f}
 \leq j!\,\varepsilon^{-j}\bigl(e^{-(\lambda-2\varepsilon u)\cdot f}
            +e^{-(\lambda+2\varepsilon u)\cdot f}\bigr),
\]
using \(e^{2\varepsilon|s|}\leq e^{2\varepsilon s}+e^{-2\varepsilon s}\).
Call the right side \(G_j\); it is \(\nu\)-integrable. Let
\(\zeta(t)=Z_\nu(\lambda+tu)=\int e^{-\lambda\cdot f}e^{-ts}\,d\nu\). For
\(|t|<\varepsilon\) and small \(h\neq0\), the mean value theorem writes the
difference quotient of the integrand as \(-s\,e^{-\tau s}e^{-\lambda\cdot f}\) for
some \(\tau\) between \(t\) and \(t+h\); its absolute value is at most \(G_1\),
and it converges to \(-s\,e^{-(\lambda+tu)\cdot f}\) as \(h\to0\). Dominated
convergence gives \(\zeta'(t)=-\int s\,e^{-(\lambda+tu)\cdot f}d\nu\) for
\(|t|<\varepsilon\). The same argument applied to \(\zeta'\), with \(G_2\) as
dominating function, gives \(\zeta''(t)=\int s^2e^{-(\lambda+tu)\cdot f}d\nu\).
Dividing by \(\zeta(0)=Z_\nu(\lambda)\),
\[
 (\ln\zeta)'(0)=\frac{\zeta'(0)}{\zeta(0)}=-E_{Q_\lambda}[s],\qquad
 (\ln\zeta)''(0)=\frac{\zeta''(0)}{\zeta(0)}-\Bigl(\frac{\zeta'(0)}{\zeta(0)}\Bigr)^2
   =\operatorname{Var}_{Q_\lambda}(s).
\]
With \(u=e_k\) the bound \(G_2\) shows \(f_k^2q_\lambda\) integrable, so
\(f_k\in L^2(Q_\lambda)\subseteq L^1(Q_\lambda)\), and
\(|f_jf_k|\leq\tfrac12(f_j^2+f_k^2)\) gives integrability of products. Taking
\(u=e_r\) gives (2-11) and (2-14).

(3) If \(\operatorname{Var}_{Q_\lambda}(u\cdot f)=0\), then \(u\cdot f\) equals
its mean \(Q_\lambda\)-almost surely. A \(Q_\lambda\)-null set \(N\) satisfies
\(\int_Nq_\lambda\,d\nu=0\) with \(q_\lambda>0\), so \(\nu(N)=0\). Thus
\(u\cdot f\) would be constant \(\nu\)-almost everywhere, which minimality
excludes for \(u\neq0\). For a segment from \(\lambda\) to \(\lambda'\neq\lambda\)
inside the open interior, put \(g(t)=\Psi_c(\lambda+t(\lambda'-\lambda))\) on a
slightly larger open interval of \(t\). By (2) at each point of the segment,
\(g\) is twice differentiable with
\(g''(t)=\operatorname{Var}_{Q_{\lambda+t(\lambda'-\lambda)}}((\lambda'-\lambda)\cdot f)>0\);
the term \(\lambda\cdot c\) is affine. So \(g'\) is strictly increasing and
\(g\) is strictly convex.

(4) If \(u\cdot f=\kappa\) \(\nu\)-almost everywhere, then
\(e^{-(\lambda+su)\cdot f}=e^{-s\kappa}e^{-\lambda\cdot f}\) almost everywhere;
hence \(Z_\nu(\lambda+su)=e^{-s\kappa}Z_\nu(\lambda)\) and the normalized
densities agree almost everywhere.
\end{proof}
\end{annexitem}

\begin{annexitem}{annex:p15-w02b-p15-043-interior}{thm:p15-w02b-p15-043-interior}{Moments of positive densities are interior}
\begin{proof}
Write \(|\cdot|\) for the Euclidean norm and \(\Sigma=E_P[(f-c)(f-c)^{\mathsf T}]\),
which exists because \(f\in L^2(P)\). For \(u\neq0\),
\(u^{\mathsf T}\Sigma u=E_P[(u\cdot f-u\cdot c)^2]\). If this were \(0\),
\(u\cdot f\) would equal \(u\cdot c\) outside a \(P\)-null set; as \(p>0\)
almost everywhere, a \(P\)-null set is \(\nu\)-null, contradicting minimality.
So \(\Sigma\) is positive definite.

For \(R>0\) let \(h_R=f\,\mathbf 1_{\{|f|\leq R\}}\), so \(|h_R|\leq R\), let
\(\bar h_R=E_P[h_R]\), and let \(C_R=E_P[(f-c)(h_R-\bar h_R)^{\mathsf T}]\). As
\(R\to\infty\), \(h_R\to f\) pointwise with \(|h_{R,k}|\leq|f_k|\), so
\(\bar h_R\to c\) by dominated convergence; the \((j,k)\) entry of the integrand
converges to \((f_j-c_j)(f_k-c_k)\) and is dominated by
\(|f_j-c_j|\,(|f_k|+E_P|f_k|)\), which is integrable because \(f\in L^2(P)\).
Hence \(C_R\to\Sigma\), and since the determinant is continuous, \(C=C_R\) is
invertible for some fixed \(R\). Let \(\beta\) be the operator norm of \(C^{-1}\).

Let \(v\in\mathbb R^m\) with \(|v|<1/(2R\beta)\) and \(w=C^{-1}v\), so
\(|w|<1/(2R)\). Define \(p_v=p\,(1+w\cdot(h_R-\bar h_R))\). Since
\(|h_R-\bar h_R|\leq2R\), the factor lies in \((0,2)\), so \(p_v\geq0\) and
\(p_v\leq2p\). Its integral is \(1+w\cdot(E_P[h_R]-\bar h_R)=1\). Thus \(p_v\)
is the density of a probability \(P_v\ll\nu\), each \(f_k\) is
\(P_v\)-integrable, and, because \(E_P[c\,(h_R-\bar h_R)^{\mathsf T}]=0\),
\[
 E_{P_v}[f]=c+E_P\bigl[(f-c)(h_R-\bar h_R)^{\mathsf T}\bigr]w=c+Cw=c+v .
\]
Hence the open ball of radius \(1/(2R\beta)\) about \(c\) lies in
\(\mathcal M_\nu\).
\end{proof}
\end{annexitem}

\begin{annexitem}{annex:p15-w02b-p15-043-existence}{thm:p15-w02b-p15-043-existence}{Existence and uniqueness of multipliers}
\begin{proof}
Write \(\mu(\lambda)=E_{Q_\lambda}[f]\). Since \(\Lambda_\nu\) is open, each of
its points is interior, so \cref{thm:p15-w02b-p15-043-partition} applies at
every point.

\emph{Into the interior.} For \(\lambda\in\Lambda_\nu\), \(Q_\lambda\) has a
positive density and \(f\in L^2(Q_\lambda)\), and
\(Q_\lambda\in\mathcal F_\nu(\mu(\lambda))\). By
\cref{thm:p15-w02b-p15-043-interior}, \(\mu(\lambda)\) is interior to
\(\mathcal M_\nu\).

\emph{Injectivity.} Suppose \(\mu(\lambda)=\mu(\lambda')=c\) with
\(\lambda\neq\lambda'\). The segment between them lies in the convex set
\(\Lambda_\nu\). For \(g(t)=\Psi_c(\lambda+t(\lambda'-\lambda))\),
\cref{thm:p15-w02b-p15-043-partition}(2) gives
\(g'(t)=(\lambda'-\lambda)\cdot(c-\mu(\lambda+t(\lambda'-\lambda)))\), so
\(g'(0)=g'(1)=0\). This contradicts the strict increase of \(g'\) from part (3).

\emph{Surjectivity.} Let \(c\) be interior to \(\mathcal M_\nu\). Regard
\(\Psi_c\) as a function \(\mathbb R^m\to(-\infty,+\infty]\). It is convex,
because \(\Lambda_\nu\) is convex and \(\ln Z_\nu\) is convex on it, and lower
semicontinuous, because \(Z_\nu\) is and the remaining term is linear. Fix
\(\lambda^*\in\Lambda_\nu\).

Step 1. Let \(|u|=1\) and \(A=\{x:u\cdot(f(x)-c)<0\}\). If \(\nu(A)=0\), then
\(u\cdot(f-c)\geq0\) almost everywhere, so \(u\cdot(E_Pf-c)\geq0\) for every
\(P\) in the definition of \(\mathcal M_\nu\). But \(c-\delta u\in\mathcal M_\nu\)
for small \(\delta>0\), and \(u\cdot(-\delta u)<0\). Hence \(\nu(A)>0\), and by
continuity from below \(\nu(A_\eta)>0\) for some \(\eta>0\), where
\(A_\eta=\{u\cdot(f-c)<-\eta\}\). For \(t>0\),
\[
 \Psi_c(\lambda^*+tu)=\ln\int e^{-\lambda^*\cdot f}e^{-tu\cdot(f-c)}\,d\nu+\lambda^*\cdot c
 \geq t\eta+\ln\int_{A_\eta}e^{-\lambda^*\cdot f}\,d\nu+\lambda^*\cdot c ,
\]
and the last integral lies in \((0,Z_\nu(\lambda^*)]\). So
\(\Psi_c(\lambda^*+tu)\to\infty\) as \(t\to\infty\).

Step 2. The set \(L=\{\lambda:\Psi_c(\lambda)\leq\Psi_c(\lambda^*)\}\) is bounded.
Otherwise choose \(\lambda_j\in L\) with \(|\lambda_j-\lambda^*|\to\infty\) and,
passing to a subsequence, \(u_j=(\lambda_j-\lambda^*)/|\lambda_j-\lambda^*|\to u\)
with \(|u|=1\). Fix \(t>0\). For large \(j\), \(\lambda^*+tu_j\) lies on the
segment from \(\lambda^*\) to \(\lambda_j\), so by convexity
\(\Psi_c(\lambda^*+tu_j)\leq\Psi_c(\lambda^*)\). Lower semicontinuity gives
\(\Psi_c(\lambda^*+tu)\leq\Psi_c(\lambda^*)\) for every \(t>0\), contradicting
Step 1.

Step 3. \(L\) is closed by lower semicontinuity, bounded and nonempty. A minimizing
sequence in \(L\) has a convergent subsequence, and lower semicontinuity makes its
limit \(\hat\lambda\) a minimizer of \(\Psi_c\) on \(L\), hence on \(\mathbb R^m\).
Since \(\Psi_c(\hat\lambda)<\infty\), \(\hat\lambda\in\Lambda_\nu\). For each
\(k\), \(t\mapsto\Psi_c(\hat\lambda+te_k)\) is differentiable at \(t=0\) and
minimal there, so its derivative \(c_k-E_{Q_{\hat\lambda}}[f_k]\) is \(0\).
Thus \(\mu(\hat\lambda)=c\).

\emph{Consequences.} (1) By bijectivity the solution \(\lambda(c)\) is unique.
It makes \(Q_{\lambda(c)}\) feasible, so
\cref{thm:p15-w02b-p15-043-optimality} makes it the unique maximizer with
\(v_\nu(c)=\Psi_c(\lambda(c))\). By \cref{thm:p15-w02b-p15-043-gibbs}(3),
\(\Psi_c(\lambda)\geq v_\nu(c)\) for every \(\lambda\in\Lambda_\nu\), so
\(\lambda(c)\) minimizes \(\Psi_c\). (2) A solution \(\lambda\in\Lambda_\nu\)
would give \(c=\mu(\lambda)\), interior by the first part. A maximizer of the form
\(Q_\lambda\) is feasible, so it would also give \(\mu(\lambda)=c\). (3) If
\(c\notin\mathcal M_\nu\), the feasible set is empty by definition, and a solution
would make \(Q_\lambda\) feasible.
\end{proof}
\end{annexitem}

\begin{annexitem}{annex:p15-w02b-p15-043-finite}{thm:p15-w02b-p15-043-finite}{The finite case, including boundary values}
\begin{proof}
(1) Every probability on \(\mathcal X\) is absolutely continuous with respect
to \(\nu\), since \(w>0\). Its moment vector \(\sum_xP\{x\}f(x)\) is a convex
combination of the points \(f(x)\), and every such convex combination arises
from the probability with those weights. Also
\(Z_\nu(\lambda)=\sum_xw(x)e^{-\lambda\cdot f(x)}<\infty\) for every
\(\lambda\).

(2) Identify probabilities with vectors in the simplex
\(\Delta\subset\mathbb R^{\mathcal X}\). The feasible set is the intersection of
\(\Delta\) with the closed set \(\{\sum_xP\{x\}f(x)=c\}\), hence compact, and
nonempty for \(c\in\mathcal M_\nu\). The entropy
\(-\sum_xP\{x\}\ln(P\{x\}/w(x))\) is a finite sum of functions continuous on
\([0,1]\), with value \(0\) at \(P\{x\}=0\). It attains a maximum on the
feasible set, every value is finite, and uniqueness follows from
\cref{thm:p15-w02b-p15-043-uniqueness}. The bound follows from
\cref{thm:p15-w02b-p15-043-no-constraint}, since every feasible \(P\) is a
competitor in the unconstrained problem.

(3) Let \(P^*\) be the maximizer. Its support lies in \(S(c)\) because \(P^*\)
is feasible. Suppose some feasible \(P\) has \(P\{x_0\}>0=P^*\{x_0\}\). For
\(t\in[0,1]\), \(P_t=(1-t)P^*+tP\) is feasible. Write \(a_x(t)=P_t\{x\}\) and
\(\eta_x(a)=-a\ln(a/w(x))\). For \(x\) with \(P^*\{x\}>0\), \(a_x(t)\geq
P^*\{x\}/2>0\) for \(t\leq\tfrac12\), so \(\eta_x(a_x(t))\) is continuously
differentiable on \([0,\tfrac12]\); hence there is \(C\) with
\(\eta_x(a_x(t))-\eta_x(a_x(0))\geq-Ct\) for all these \(x\) and
\(t\in[0,\tfrac12]\). For \(x\) with \(P^*\{x\}=0\), \(a_x(t)=tP\{x\}\) and
\(\eta_x(tP\{x\})\geq0\) once \(tP\{x\}\leq w(x)\). Therefore, for small \(t>0\),
\[
 H_\nu(P_t)-H_\nu(P^*)\geq t\Bigl(P\{x_0\}\ln\frac{w(x_0)}{tP\{x_0\}}-C\Bigr),
\]
which is positive once \(\ln(1/t)\) is large. This contradicts maximality.
So the support of \(P^*\) is \(S(c)\).

(4) Assume minimality; \(\Lambda_\nu=\mathbb R^m\) is open and nonempty. If
\(c\) is interior, \cref{thm:p15-w02b-p15-043-existence} gives a solution and
the maximizer \(Q_{\lambda(c)}\), whose density is positive at every point. If
the multiplier equations have a solution, the same theorem shows that \(c\) is
interior. If the maximizer gives positive probability to every point, its
density relative to \(\nu\) is positive and \(f\) is bounded, so
\cref{thm:p15-w02b-p15-043-interior} makes \(c\) interior.
\end{proof}
\end{annexitem}

\begin{annexitem}{annex:p15-w02b-p15-043-failures}{thm:p15-w02b-p15-043-failures}{Nonexistence and nonattainment}
\begin{proof}
(1) Feasibility means \(P\{1\}+2P\{2\}=2\) and \(P\{0\}+P\{1\}+P\{2\}=1\).
Subtracting twice the second from the first gives \(-2P\{0\}-P\{1\}=0\), so
\(P=\delta_2\), whose entropy is \(0\). For every real \(\lambda\),
\[
 2-E_{Q_\lambda}[x]=2-\frac{e^{-\lambda}+2e^{-2\lambda}}{1+e^{-\lambda}+e^{-2\lambda}}
 =\frac{2+e^{-\lambda}}{1+e^{-\lambda}+e^{-2\lambda}}>0 ,
\]
so the multiplier equation has no solution.

(2) If \(\lambda\geq0\), \(\int_{-\infty}^0e^{-\lambda x}dx\geq\int_{-\infty}^01\,dx=\infty\);
if \(\lambda<0\), \(\int_0^\infty e^{-\lambda x}dx=\infty\) likewise. So
\(\Lambda_\nu=\varnothing\). For \(L>0\) the uniform density \(1/(2L)\) on
\([c-L,c+L]\) is feasible with mean \(c\) and entropy \(\ln(2L)\), which tends to
\(\infty\). So \(v_\nu(c)=+\infty\), and any feasible distribution of finite
entropy is exceeded.

(3) \emph{Domain.} For \(\lambda\geq0\),
\(Z_\nu(\lambda)\leq\int_0^\infty(1+x)^{-3}dx=\tfrac12\). For \(\lambda<0\) and
\(x\geq1\), \(e^{|\lambda|x}\geq|\lambda|^4x^4/24\) and \((1+x)^{-3}\geq(2x)^{-3}\),
so the integrand is at least \(|\lambda|^4x/192\) and \(Z_\nu(\lambda)=\infty\).
So \(\Lambda_\nu=[0,\infty)\), which contains \(0\) but no negative number.

\emph{Moment set.} Since the density of \(\nu\) is positive, \(\nu\) and
Lebesgue measure on \([0,\infty)\) have the same null sets. A probability
\(P\ll\nu\) gives \(\{0\}\) probability \(0\), so \(x>0\) \(P\)-almost surely
and \(E_P[x]>0\) when finite. Conversely the uniform law on \([0,2c]\) has
mean \(c\). So \(\mathcal M_\nu=(0,\infty)\), and every \(c>0\) is interior.

\emph{Means of the family.} \(Q_0\) has Lebesgue density \(2(1+x)^{-3}\) and
mean \(2\int_0^\infty\bigl((1+x)^{-2}-(1+x)^{-3}\bigr)dx=2(1-\tfrac12)=1\).
For \(\lambda>0\), \((x-1)(e^{-\lambda x}-e^{-\lambda})\leq0\) for every
\(x\geq0\), since the two factors have opposite signs. Taking
\(Q_0\)-expectations and using \(E_{Q_0}[x]=1\),
\[
 E_{Q_0}[xe^{-\lambda x}]-E_{Q_0}[e^{-\lambda x}]\leq e^{-\lambda}\bigl(E_{Q_0}[x]-1\bigr)=0 .
\]
Since \(q_\lambda/q_0=e^{-\lambda x}/E_{Q_0}[e^{-\lambda x}]\), this says
\(E_{Q_\lambda}[x]\leq1\). Hence for \(c>1\) no \(\lambda\in\Lambda_\nu\) gives
mean \(c\); in particular the multiplier equation has no solution at any
\(\lambda>0\).

\emph{Value.} By \cref{thm:p15-w02b-p15-043-gibbs} with \(\lambda=0\), every
feasible \(P\) satisfies \(H_\nu(P)=-\ln2-D(P\Vert Q_0)\). As \(P\) has mean
\(c\neq1\), \(P\neq Q_0\) and \(H_\nu(P)<-\ln2\). For \(a>c\), let \(R_a\) be
\(Q_0\) conditioned on \([a,a+1]\), with mean \(m_a\in[a,a+1]\), let
\(\varepsilon_a=(c-1)/(m_a-1)\in(0,1)\), and let
\(P_a=(1-\varepsilon_a)Q_0+\varepsilon_aR_a\). Then \(P_a\ll\nu\) and its mean is
\(1+\varepsilon_a(m_a-1)=c\). Its density relative to \(Q_0\) is
\(1-\varepsilon_a+\varepsilon_ar\), where
\(r=\mathbf 1_{[a,a+1]}/Q_0[a,a+1]\). Since \(t\mapsto t\ln t\) is convex,
\((1-\varepsilon_a+\varepsilon_ar)\ln(1-\varepsilon_a+\varepsilon_ar)\leq
\varepsilon_ar\ln r\), and integrating against \(Q_0\),
\[
 D(P_a\Vert Q_0)\leq\varepsilon_a\ln\frac1{Q_0[a,a+1]}
 \leq\frac{c-1}{a-1}\bigl(3\ln(2+a)-\ln2\bigr),
\]
because \(Q_0[a,a+1]=2\int_a^{a+1}(1+x)^{-3}dx\geq2(2+a)^{-3}\). The right side
tends to \(0\) as \(a\to\infty\). So \(H_\nu(P_a)\to-\ln2\) and
\(v_\nu(c)=-\ln2=\ln Z_\nu(0)=\Psi_c(0)\). By
\cref{thm:p15-w02b-p15-043-gibbs}(3), \(\Psi_c\geq v_\nu(c)=\Psi_c(0)\) on
\(\Lambda_\nu\), so the infimum is \(\Psi_c(0)\). The value is not attained.

(4) By (1) of \cref{thm:p15-w02b-p15-043-finite}, \(\mathcal M_\nu=[0,2]\), so
\(3\notin\mathcal M_\nu\), the feasible set is empty and
\(v_\nu(3)=-\infty\). Moreover
\(\Psi_3(\lambda)=\ln(1+e^{-\lambda}+e^{-2\lambda})+3\lambda
 =\ln(1+e^{\lambda}+e^{2\lambda})+\lambda\), which tends to \(-\infty\) as
\(\lambda\to-\infty\).
\end{proof}
\end{annexitem}

\begin{annexitem}{annex:p15-w02b-p15-043-transport}{thm:p15-w02b-p15-043-transport}{Reparametrization and rescaling of the reference measure}
\begin{proof}
(1) For measurable \(g\geq0\) on \(\mathcal Y\),
\(\int g\,d(T_*\nu)=\int g\circ T\,d\nu\): for indicators this is the
definition \(T_*\nu(B)=\nu(T^{-1}B)\); it extends to nonnegative simple
functions by linearity and to nonnegative measurable functions by monotone
convergence applied to increasing simple approximations. The measure
\(T_*\nu\) is \(\sigma\)-finite, since \(T\) maps a cover of \(\mathcal X\) by
sets of finite \(\nu\)-measure to such a cover of \(\mathcal Y\), the image sets
being measurable because \(T^{-1}\) is measurable. If \(P\ll\nu\) has density
\(p\), then for \(B\in\mathcal B\)
\[
 T_*P(B)=\int\mathbf 1_B\circ T\,p\,d\nu
        =\int(\mathbf 1_B\cdot p\circ T^{-1})\circ T\,d\nu
        =\int_Bp\circ T^{-1}\,d(T_*\nu),
\]
so \(T_*P\ll T_*\nu\) with density \(\tilde p=p\circ T^{-1}\). The same formula
applied to \(|f_k\circ T^{-1}|\tilde p\) and to \((f_k\circ T^{-1})^\pm\tilde p\)
shows \(f_k\circ T^{-1}\in L^1(T_*P)\) with
\(E_{T_*P}[f_k\circ T^{-1}]=E_P[f_k]\). Applied to \((\tilde p\ln\tilde p)^\pm
=(p\ln p)^\pm\circ T^{-1}\), it shows that the two entropies are defined together
and are equal. The inverse correspondence is given by \(T^{-1}\), since
\((T^{-1})_*T_*\nu=\nu\). Feasibility and entropy correspond, so maximizers do.

(2) A density of \(P\) relative to \(a\nu\) is \(p/a\), and
\((p/a)\ln(p/a)\,a=p\ln p-p\ln a\), where \(p\ln a\) is integrable with integral
\(\ln a\). So the entropies are defined together and
\(H_{a\nu}(P)=H_\nu(P)+\ln a\). The feasible sets coincide, because \(\nu\) and
\(a\nu\) have the same null sets. So the maximizers coincide.

(3) \(T\) is a continuous bijection of \((0,\infty)\) with continuous inverse
\(u\mapsto\sqrt{2u}\), hence bimeasurable. For \(0\leq\alpha<\beta\),
\(T^{-1}(\alpha,\beta]=(\sqrt{2\alpha},\sqrt{2\beta}]\), and
\[
 \int_{\sqrt{2\alpha}}^{\sqrt{2\beta}}x\,dx=\beta-\alpha,\qquad
 \int_{\sqrt{2\alpha}}^{\sqrt{2\beta}}dx=\sqrt{2\beta}-\sqrt{2\alpha}
 =\int_\alpha^\beta(2u)^{-1/2}\,du .
\]
Fix \(N>0\). On \((0,N]\) the image of \(x\,dx\) and Lebesgue measure are finite
measures with the same total mass that agree on the intervals \((\alpha,\beta]\)
with \(0\leq\alpha<\beta\leq N\). These intervals with the empty set form a
\(\pi\)-system generating the Borel sets of \((0,N]\). After division by the
common total mass, \cref{thm:p15-measure-uniqueness} makes the two measures equal
on \((0,N]\). The same argument applies to the image of \(dx\) and
\((2u)^{-1/2}du\). Letting \(N\to\infty\) through the integers and using countable
additivity gives equality on \((0,\infty)\).
\end{proof}
\end{annexitem}

\begin{annexitem}{annex:p15-w02b-p15-043-shannon}{thm:p15-w02b-p15-043-shannon}{The uncertainty measure from Jaynes's Appendix~A conditions}
\begin{proof}
\emph{Multiplicativity.} Apply (3) to the uniform vector of length \(mn\),
grouped into \(m\) consecutive groups of \(n\). Each group has probability
\(1/m\) and conditional vector \((1/n,\ldots,1/n)\), so
\[
 A(mn)=A(m)+\sum_{j=1}^m\tfrac1mA(n)=A(m)+A(n),
\]
which is Jaynes's (A-3) \cite[p.~630]{Jaynes1957InformationTheory}. With
\(m=n=1\), \(A(1)=0\); by induction \(A(n^t)=tA(n)\) for integers
\(t\geq1\).

\emph{The logarithm.} By (2), \(A(2)>A(1)=0\); put \(K=A(2)/\ln2>0\). Let
\(n\geq2\) and \(t\geq1\), and choose the integer \(r\geq0\) with
\(2^r\leq n^t<2^{r+1}\). Monotonicity gives
\(rA(2)\leq tA(n)\leq(r+1)A(2)\), and the logarithm gives
\(r\ln2\leq t\ln n<(r+1)\ln2\). So both \(A(n)/A(2)\) and \(\ln n/\ln2\) lie in
\([r/t,(r+1)/t]\) and differ by at most \(1/t\). As \(t\) is arbitrary,
\(A(n)=K\ln n\), which is (A-4); this also holds for \(n=1\). Jaynes refers to
Shannon for the uniqueness of this solution
\cite[p.~630]{Jaynes1957InformationTheory}; the argument above supplies it from
the monotonicity condition.

\emph{Rational vectors.} Let \(n_1,\ldots,n_k\) be positive integers,
\(N=\sum_in_i\) and \(p_i=n_i/N\). Group the uniform vector of length \(N\)
into consecutive groups of sizes \(n_1,\ldots,n_k\). The group probabilities are
\(p_i>0\) and the conditional vectors are uniform of lengths \(n_i\), so (3)
gives (A-2): \(A(N)=H(p_1,\ldots,p_k)+\sum_ip_iA(n_i)\). Hence
\[
 H(p_1,\ldots,p_k)=K\ln N-K\sum_ip_i\ln n_i=-K\sum_ip_i\ln p_i ,
\]
which is (A-5).

\emph{All vectors.} For fixed \(k\), probability vectors with positive rational
entries are dense in the simplex: a vector \(p\) is the limit of
\((1-1/j)p+(1/j)(1/k,\ldots,1/k)\), which has positive entries, and a vector with
positive entries is approximated by choosing rationals near its first \(k-1\)
entries and letting the last entry be one minus their sum. Both \(H\), by (1), and
\(-K\sum_ip_i\ln p_i\), with \(0\ln0=0\), are continuous on the simplex and agree
on a dense subset, so they agree everywhere.
\end{proof}
\end{annexitem}

\begin{annexitem}{annex:p15-w02b-p15-043-power-family}{thm:p15-w02b-p15-043-power-family}{Maximum entropy against the reference measures \(x^n\,dx\)}
\begin{proof}
For \(z>0\) let \(\Gamma(z)=\int_0^\infty y^{z-1}e^{-y}dy\), which is finite;
integration by parts gives \(\Gamma(z+1)=z\Gamma(z)\).

(1) If \(\lambda\leq0\), the integrand \(x^ne^{-\lambda x}\) is at least \(x^n\)
on \([1,\infty)\), whose integral diverges because \(n>-1\). If \(\lambda>0\),
the substitution \(y=\lambda x\) gives
\(\int_0^\infty x^{n+j}e^{-\lambda x}dx=\Gamma(n+j+1)\lambda^{-(n+j+1)}\) for
\(j=0,1,2\). So \(\Lambda_{\nu_n}=(0,\infty)\),
\(Z(\lambda)=\Gamma(n+1)\lambda^{-(n+1)}\), and
\[
 E_{Q_\lambda}[x]=\frac{\Gamma(n+2)}{\Gamma(n+1)\lambda}=\frac{n+1}\lambda,
 \qquad
 E_{Q_\lambda}[x^2]=\frac{(n+2)(n+1)}{\lambda^2}.
\]
The equation \((n+1)/\lambda=c\) has the single solution \(\lambda=(n+1)/c\),
at which \(Q_\lambda\) is feasible. \Cref{thm:p15-w02b-p15-043-optimality} makes
it the unique maximizer with value
\(\ln Z(\lambda)+\lambda c=\ln\Gamma(n+1)-(n+1)\ln\lambda+n+1\). Its Lebesgue
density is \(x^nq_\lambda(x)=\lambda^{n+1}x^ne^{-\lambda x}/\Gamma(n+1)\), and
\(E[x^2]=(n+2)(n+1)c^2/(n+1)^2=(n+2)c^2/(n+1)\).

(2) Divergence for \(\lambda\leq0\) is as in (1). For \(\lambda>0\) the
substitution \(y=\lambda x^2\), \(dx=\tfrac12(\lambda y)^{-1/2}dy\), gives
\[
 I_j\defeq\int_0^\infty x^{n+j}e^{-\lambda x^2}dx
 =\tfrac12\lambda^{-(n+j+1)/2}\,\Gamma\bigl(\tfrac{n+j+1}2\bigr)
 \qquad(j=0,1,2).
\]
So \(Z(\lambda)=I_0\), \(E_{Q_\lambda}[x^2]=I_2/I_0=\Gamma(\tfrac{n+3}2)/
(\lambda\Gamma(\tfrac{n+1}2))=(n+1)/(2\lambda)\), and
\(E_{Q_\lambda}[x]=I_1/I_0=\lambda^{-1/2}\Gamma(\tfrac{n+2}2)/\Gamma(\tfrac{n+1}2)\).
The equation \((n+1)/(2\lambda)=s\) has the single solution
\(\lambda=(n+1)/(2s)\); \(Q_\lambda\) is feasible, and
\cref{thm:p15-w02b-p15-043-optimality} makes it the unique maximizer with value
\(\ln Z(\lambda)+\lambda s\). Substituting \(\lambda\) gives the stated
\(E[x]\), and the Lebesgue density is \(x^n e^{-\lambda x^2}/I_0\).

\emph{Agreement with the source.} For integer \(n\geq0\), Jaynes's partition
function \(n!/(\epsilon\lambda^{n+1})\) is \(Z(\lambda)/\epsilon\) because
\(\Gamma(n+1)=n!\), and his (3-10) is the formula for \(E[x^2]\)
\cite[p.~625]{Jaynes1957InformationTheory}. Writing \(k!\) for \(\Gamma(k+1)\) at
half-integers as he does, \((n/2)!=\Gamma(\tfrac{n+2}2)\) and
\([\tfrac12(n+1)]!=\tfrac{n+1}2\Gamma(\tfrac{n+1}2)\), so his coefficient
\(\bigl(\tfrac{n+1}2\bigr)^{1/2}(\tfrac12n)!/[\tfrac12(n+1)]!\) in (3-11) equals
\(\bigl(\tfrac2{n+1}\bigr)^{1/2}\Gamma(\tfrac{n+2}2)/\Gamma(\tfrac{n+1}2)\), the
coefficient of \(s^{1/2}\) above. The constant \(1/\epsilon\) leaves the
maximizers unchanged by \cref{thm:p15-w02b-p15-043-transport}(2).
\end{proof}
\end{annexitem}
